\subsection{Sequential Monte Carlo and particle filters}
\label{subsec:smc}

The earlier methods in this unit target one fixed distribution.
State-space models change the problem: after observing $y_1,\dots,y_t$ we want the current posterior for the latent state, and at time $t+1$ that target changes because a new observation arrives.
Sequential Monte Carlo (SMC) is built for exactly this update cycle.
At each time step it represents the current filtering law by particles, reweights those particles using the new observation, and resamples when the weights have become too uneven.

The augmentation is therefore different from the momentum and temperature variables in the previous chapters.
Here the auxiliary state is a whole interacting population, together with weights and eventually ancestry variables, that tracks an evolving sequence of posterior distributions.
The particles propagate through the latent dynamics, the weights incorporate the newest observation, and resampling prevents the approximation from collapsing onto a single lineage.

\subsubsection{State-space models and filtering recursions}
\label{subsubsec:smc-ssm}

Let $x_t\in\R^d$ denote a latent state and $y_t\in\R^m$ the corresponding observation.
A state-space model is specified by an initial law $p_0$, a transition density $f_\theta$, and an observation density $g_\theta$:
\begin{equation}
  x_0 \sim p_0,
  \qquad
  x_t \mid x_{t-1} \sim f_\theta(\cdot \mid x_{t-1}),
  \qquad
  y_t \mid x_t \sim g_\theta(\cdot \mid x_t).
  \label{eq:smc-ssm}
\end{equation}
The latent process $(x_t)_{t\ge 0}$ is Markov, and the observations are conditionally independent given the states.
For data $y_{1:T}$, the four recurring inferential tasks are
\[
  \text{filtering: } p_\theta(x_t \mid y_{1:t}),
  \qquad
  \text{prediction: } p_\theta(x_{t+k} \mid y_{1:t}),
\]
\[
  \text{smoothing: } p_\theta(x_{0:t} \mid y_{1:T}),
  \qquad
  \text{parameter learning: infer } \theta \text{ from } y_{1:T}.
\]
Filtering is the central online problem, because once we can update $p_\theta(x_t\mid y_{1:t})$ sequentially, prediction and likelihood evaluation follow from the same recursion, while smoothing and parameter learning can be built on top of those filtered quantities.

\begin{proposition}[Prediction and update]
\label{prop:smc-filtering-recursion}
For the state-space model \eqref{eq:smc-ssm}, the predictive and filtering laws satisfy
\begin{equation}
  p_\theta(x_t \mid y_{1:t-1})
  =
  \int f_\theta(x_t \mid x_{t-1})\,p_\theta(x_{t-1} \mid y_{1:t-1})\,dx_{t-1},
  \label{eq:smc-predict}
\end{equation}
and
\begin{equation}
  p_\theta(x_t \mid y_{1:t})
  =
  \frac{g_\theta(y_t \mid x_t)\,p_\theta(x_t \mid y_{1:t-1})}{z_t},
  \qquad
  z_t \coloneqq p_\theta(y_t \mid y_{1:t-1})
  =
  \int g_\theta(y_t \mid x_t)\,p_\theta(x_t \mid y_{1:t-1})\,dx_t.
  \label{eq:smc-update}
\end{equation}
\end{proposition}
\begin{proof}
Using the Markov property of the latent chain,
\[
  p_\theta(x_t \mid x_{t-1},y_{1:t-1}) = p_\theta(x_t \mid x_{t-1}) = f_\theta(x_t \mid x_{t-1}).
\]
Therefore
\[
  p_\theta(x_t \mid y_{1:t-1})
  =
  \int p_\theta(x_t \mid x_{t-1},y_{1:t-1})\,p_\theta(x_{t-1} \mid y_{1:t-1})\,dx_{t-1},
\]
which gives \eqref{eq:smc-predict}.
For the update, Bayes's rule and conditional independence of $y_t$ from $y_{1:t-1}$ given $x_t$ yield
\[
  p_\theta(x_t \mid y_{1:t})
  =
  \frac{p_\theta(y_t \mid x_t,y_{1:t-1})\,p_\theta(x_t \mid y_{1:t-1})}{p_\theta(y_t \mid y_{1:t-1})}
  =
  \frac{g_\theta(y_t \mid x_t)\,p_\theta(x_t \mid y_{1:t-1})}{p_\theta(y_t \mid y_{1:t-1})}.
\]
Integrating the numerator over $x_t$ gives the normalizing constant $z_t$ and proves \eqref{eq:smc-update}.%
\end{proof}

\subsubsection{Kalman filtering as a tractable baseline}
\label{subsubsec:smc-kalman}

The prediction-update recursion becomes especially simple in the linear-Gaussian model
\begin{equation}
  x_0 \sim \mathcal{N}(m_0,P_0),
  \qquad
  x_t = A x_{t-1} + q_t,
  \quad q_t \sim \mathcal{N}(0,Q),
  \qquad
  y_t = C x_t + r_t,
  \quad r_t \sim \mathcal{N}(0,R),
  \label{eq:smc-linear-gaussian}
\end{equation}
with independent Gaussian noises.
If
\[
  x_{t-1} \mid y_{1:t-1} \sim \mathcal{N}(m_{t-1},P_{t-1}),
\]
then \eqref{eq:smc-predict} implies
\begin{equation}
  x_t \mid y_{1:t-1} \sim \mathcal{N}(\hat m_t,\hat P_t),
  \qquad
  \hat m_t = A m_{t-1},
  \qquad
  \hat P_t = A P_{t-1} A^\top + Q.
  \label{eq:smc-kalman-predict}
\end{equation}
The update is equally explicit.
Conditional on $y_{1:t-1}$, the observation equation in \eqref{eq:smc-linear-gaussian} gives
\[
  y_t = Cx_t + r_t,
  \qquad
  r_t \sim \mathcal{N}(0,R),
\]
with $r_t$ independent of $x_t$.
Therefore the pair $(x_t,y_t)$ is jointly Gaussian with
\begin{equation}
  \begin{bmatrix}
    x_t \\
    y_t
  \end{bmatrix}
  \,\bigm|\,
  y_{1:t-1}
  \sim
  \mathcal{N}\!\left(
    \begin{bmatrix}
      \hat m_t \\
      C\hat m_t
    \end{bmatrix},
    \begin{bmatrix}
      \hat P_t & \hat P_t C^\top \\
      C\hat P_t & C\hat P_t C^\top + R
    \end{bmatrix}
  \right).
  \label{eq:smc-joint-gaussian}
\end{equation}
For a jointly Gaussian pair $(U,V)$ with block mean $(\mu_U,\mu_V)$ and block covariance
\[
  \begin{bmatrix}
    \Sigma_{UU} & \Sigma_{UV} \\
    \Sigma_{VU} & \Sigma_{VV}
  \end{bmatrix},
\]
completing the square in the joint density shows that
\[
  U \mid V=v
  \sim
  \mathcal{N}\!\left(
    \mu_U + \Sigma_{UV}\Sigma_{VV}^{-1}(v-\mu_V),
    \Sigma_{UU} - \Sigma_{UV}\Sigma_{VV}^{-1}\Sigma_{VU}
  \right).
\]
Applying this formula to \eqref{eq:smc-joint-gaussian} yields the innovation covariance
\[
  S_t \coloneqq C\hat P_t C^\top + R
\]
When $R \succ 0$, the matrix $S_t$ is invertible, and the Kalman gain is
\begin{equation}
  K_t \coloneqq \hat P_t C^\top S_t^{-1}
  =
  \hat P_t C^\top (C\hat P_t C^\top + R)^{-1},
  \label{eq:smc-kalman-gain}
\end{equation}
\begin{equation}
  m_t = \hat m_t + K_t(y_t - C\hat m_t),
  \qquad
  P_t = \hat P_t - K_t C\hat P_t = (I-K_tC)\hat P_t.
  \label{eq:smc-kalman-update}
\end{equation}
The correction term $y_t - C\hat m_t$ is the innovation: it measures how far the new observation lies from its prediction under the current state estimate, and the matrix $K_t$ determines how strongly that discrepancy should pull the posterior mean.
The important structural point is not the formula itself, but the Gaussian closure behind it:
linear dynamics and Gaussian noise keep the filtering distribution inside a tractable parametric family.
If the latent states $x_{0:T}$ were observed, parameter learning in \eqref{eq:smc-linear-gaussian} would reduce to ordinary multivariate regression and residual covariance estimation.
For example, one would estimate the dynamics by regressing $x_t$ on $x_{t-1}$,
\[
  \hat A
  =
  \left(\sum_{t=1}^T x_t x_{t-1}^\top\right)
  \left(\sum_{t=1}^T x_{t-1}x_{t-1}^\top\right)^{-1},
  \qquad
  \hat Q
  =
  \frac{1}{T}\sum_{t=1}^T
  (x_t-\hat A x_{t-1})(x_t-\hat A x_{t-1})^\top,
\]
while the observation model would give
\[
  \hat C
  =
  \left(\sum_{t=1}^T y_t x_t^\top\right)
  \left(\sum_{t=1}^T x_t x_t^\top\right)^{-1},
  \qquad
  \hat R
  =
  \frac{1}{T}\sum_{t=1}^T
  (y_t-\hat C x_t)(y_t-\hat C x_t)^\top.
\]
The hidden-state difficulty is exactly that these regressors are unobserved.

That is why Kalman filtering solves the online problem exactly, and why a backward Gaussian-smoothing pass is the natural bridge to the EM updates from \Cref{subsec:expectation-maximization}.
Running the Kalman filter forward gives the filtered Gaussian laws, and a backward Rauch--Tung--Striebel-type pass then returns smoothed means and covariances such as
\[
  m_t^{\mathrm{s}} \coloneqq \E_{\theta^{(k)}}[x_t\mid y_{1:T}],
  \qquad
  P_t^{\mathrm{s}} \coloneqq \operatorname{Var}_{\theta^{(k)}}(x_t\mid y_{1:T}),
  \qquad
  P_{t,t-1}^{\mathrm{s}} \coloneqq \operatorname{Cov}_{\theta^{(k)}}(x_t,x_{t-1}\mid y_{1:T}).
\]
The E-step for a current parameter value $\theta^{(k)}$ therefore reduces to computing the smoothed one-step and lag-one moments
\begin{equation}
  \E_{\theta^{(k)}}[x_t \mid y_{1:T}],
  \qquad
  \E_{\theta^{(k)}}[x_t x_t^\top \mid y_{1:T}],
  \qquad
  \E_{\theta^{(k)}}[x_t x_{t-1}^\top \mid y_{1:T}],
  \label{eq:smc-em-sufficient-statistics}
\end{equation}
which are the sufficient statistics of the complete-data Gaussian model.
Equivalently,
\[
  Q(\theta \mid \theta^{(k)})
  \coloneqq
  \E_{\theta^{(k)}}\!\left[
    \log p_\theta(x_{0:T},y_{1:T})
    \,\middle|\,
    y_{1:T}
  \right]
\]
depends on $\theta$ only through these expectations, so the M-step simply plugs them into the same regression and covariance formulas that would be used if the states were observed.
In that sense, EM for linear-Gaussian state-space models is just ``regression with missing states,'' where the backward smoother supplies the conditional replacements for the unobserved regressors and cross-products.
Once the model becomes nonlinear or non-Gaussian, the recursion from \Cref{prop:smc-filtering-recursion} still holds, but the filtered law no longer closes inside a finite-dimensional family.

\subsubsection{Particle approximations and weight degeneracy}
\label{subsubsec:smc-particles}

For a general state-space model, SMC is an online version of the self-normalized importance-sampling construction from \Cref{subsubsec:importance-sampling}.
Instead of approximating one fixed target once, it updates the approximation each time a new observation arrives.
This progression from sequential importance sampling to a general SMC framework is developed systematically in \citet[\S2.6.3 and Ch.~3]{liu_2002_monte_carlo}.
Concretely, it replaces the filtering distribution by a weighted empirical measure
\begin{equation}
  \pi_t^N(dx)
  \coloneqq
  \sum_{i=1}^N w_t^{(i)}\,\delta_{x_t^{(i)}}(dx)
  \approx
  p_\theta(x_t \mid y_{1:t})(dx),
  \label{eq:smc-empirical-measure}
\end{equation}
where the particle locations $x_t^{(i)}$ summarize likely states and the weights satisfy $w_t^{(i)}\ge 0$ and $\sum_i w_t^{(i)}=1$.
As in ordinary importance sampling, expectations under the filtering law are then approximated by weighted averages:
\begin{equation}
  \E_\theta[h(x_t)\mid y_{1:t}]
  \approx
  \sum_{i=1}^N w_t^{(i)} h(x_t^{(i)}).
  \label{eq:smc-expectation}
\end{equation}

At this point it is worth separating the empirical approximation problem from the specific sequential mechanics.
In SMC the particles interact through weighting and resampling, so after one update step the cloud is typically no longer i.i.d.
That is not itself a defect.
What matters is how well the empirical measure approximates the target law.
Suppose our task is to estimate expectations under \(\pi\) from a particle cloud of fixed size \(N\).
Then the right question is not ``are the particles independent?'' but ``what does dependence do to the variance of empirical averages?''
The next proposition records the basic variance calculation behind that point.

\begin{proposition}[Empirical averages for an exchangeable particle cloud]
\label{prop:particle-cloud-unbiased-variance}
Let \(X^{(1)},\dots,X^{(N)}\) be exchangeable random variables with common marginal law \(\pi\), and define the empirical measure
\[
  \hat\pi_N
  \coloneqq
  \frac1N\sum_{i=1}^N \delta_{X^{(i)}}.
\]
Then for every integrable test function \(h\),
\begin{equation}
  \E\!\left[\int h(x)\,\hat\pi_N(dx)\right]
  =
  \int h(x)\,\pi(dx).
  \label{eq:particle-cloud-unbiased}
\end{equation}
If \(h\in L^2(\pi)\), then
\begin{equation}
  \operatorname{Var}\!\left(\int h(x)\,\hat\pi_N(dx)\right)
  =
  \frac{1}{N}\operatorname{Var}_{\pi}\!\bigl(h(X^{(1)})\bigr)
  +
  \frac{N-1}{N}\operatorname{Cov}\!\bigl(h(X^{(1)}),h(X^{(2)})\bigr).
  \label{eq:particle-cloud-variance}
\end{equation}
\end{proposition}
\begin{proof}
Because each \(X^{(i)}\) has marginal law \(\pi\),
\[
  \E\!\left[\int h(x)\,\hat\pi_N(dx)\right]
  =
  \frac1N\sum_{i=1}^N \E[h(X^{(i)})]
  =
  \int h(x)\,\pi(dx),
\]
which proves \eqref{eq:particle-cloud-unbiased}.
For the variance, write
\[
  \int h(x)\,\hat\pi_N(dx)
  =
  \frac1N\sum_{i=1}^N h(X^{(i)}).
\]
Expanding the variance of that average gives
\[
  \operatorname{Var}\!\left(\frac1N\sum_{i=1}^N h(X^{(i)})\right)
  =
  \frac{1}{N^2}\sum_{i=1}^N \operatorname{Var}(h(X^{(i)}))
  +
  \frac{1}{N^2}\sum_{i\neq j}\operatorname{Cov}(h(X^{(i)}),h(X^{(j)})).
\]
Exchangeability makes all variance terms equal and all off-diagonal covariance terms equal, which yields \eqref{eq:particle-cloud-variance}.%
\end{proof}

For i.i.d.\ particles, the covariance term in \eqref{eq:particle-cloud-variance} vanishes.
If the dependence is negative for a class of important test functions \(h\), then the coupled cloud can approximate \(\pi\) more accurately than an i.i.d.\ cloud of the same size.
So from the viewpoint of empirical approximation, independence is sufficient but not sacred.
This matters already in SMC, where reweighting and resampling create interaction automatically and the real question is whether the resulting empirical measure is good.
Later, in \Cref{subsubsec:diffusion-diverse-sampling}, we will see a diffusion-model example where one specifies a joint law on the particle cloud directly in order to improve finite-particle approximation \citep[\S3]{corso_xu_de_bortoli_barzilay_jaakkola_2023_particle_guidance}.

If the weighted particles at time $t-1$ approximate $p_\theta(x_{t-1}\mid y_{1:t-1})$, the most direct sequential importance step uses the transition density itself as the proposal.
We first propagate
\begin{equation}
  x_t^{*,(i)} \sim f_\theta(\cdot \mid x_{t-1}^{(i)}),
  \qquad i=1,\dots,N,
  \label{eq:smc-propagate}
\end{equation}
which approximates the predictive law \eqref{eq:smc-predict}, and then reweight with the new likelihood factor:
\begin{equation}
  \tilde w_t^{(i)}
  =
  w_{t-1}^{(i)}\,g_\theta(y_t \mid x_t^{*,(i)}),
  \qquad
  w_t^{(i)}
  =
  \frac{\tilde w_t^{(i)}}{\sum_{j=1}^N \tilde w_t^{(j)}}.
  \label{eq:smc-weight-update}
\end{equation}
Equation \eqref{eq:smc-weight-update} is Bayes's rule from \eqref{eq:smc-update} written in sequential importance-sampling form, with exactly the same normalize-the-weights step as in \Cref{subsubsec:importance-sampling}.
Liu presents the same recursive-weight viewpoint in \citet[\S2.6.3]{liu_2002_monte_carlo}.
In implementation one usually evaluates the unnormalized weights on the log scale, as in \Cref{subsubsec:working-on-log-scale}, since likelihood values can vary by many orders of magnitude across particles.
The operational cycle is therefore simple: start from the weighted empirical law at time $t-1$, propagate each particle through the transition model, multiply by the new likelihood factor, renormalize, and resample only when the resulting weights become too uneven.

A standard diagnostic is the same effective-sample-size proxy from \Cref{subsubsec:proposal-quality-and-ess}, now applied at each time step.
Its role is to quantify the weight degeneracy of pure sequential importance sampling.
After many observations, most particles receive tiny weights, so the approximation is carried by only a few surviving lineages.
\begin{equation}
  \mathrm{ESS}_t
  \coloneqq
  \frac{1}{\sum_{i=1}^N (w_t^{(i)})^2}.
  \label{eq:smc-ess}
\end{equation}
This is the same weight-based ESS proxy used in importance sampling, not the autocorrelation ESS from \Cref{subsubsec:ess} for a single correlated Markov chain.
This equals $N$ when the weights are uniform and decreases toward $1$ when a single particle dominates.
Here the interpretation is even more concrete than in one-shot importance sampling:
a small ESS means the particle system has effectively forgotten most of its diversity and needs help before the next observation arrives.

\subsubsection{Resampling and the bootstrap filter}
\label{subsubsec:smc-resampling}

Resampling combats degeneracy by replacing a highly uneven weighted sample with an equally weighted sample drawn from the current empirical measure.
Given weighted particles $\{(x_t^{*,(i)},w_t^{(i)})\}_{i=1}^N$, draw ancestor indices
\[
  A_t^{(1)},\dots,A_t^{(N)} \in \{1,\dots,N\}
\]
with probabilities $\mathbb{P}(A_t^{(j)}=i)=w_t^{(i)}$, then set
\[
  x_t^{(j)} = x_t^{*,(A_t^{(j)})},
  \qquad
  w_t^{(j)} = \frac{1}{N}.
\]
The new system duplicates high-weight particles and discards low-weight ones, so future propagation starts from a healthier empirical approximation.
If the ancestor indices are retained, then the algorithm is genuinely working on an augmented state:
the current particle locations approximate the filtering law, while the recorded lineages provide the bookkeeping needed for later smoothing over past states.
Equivalently, one can formulate the method on path space, with particles carrying weighted partial trajectories $x_{0:t}^{(i)}$ rather than only their current endpoints; this is the partial-sample viewpoint emphasized by \citet[\S\S3.4.3--3.4.4]{liu_2002_monte_carlo}.

\begin{proposition}[Resampling preserves weighted averages in expectation]
\label{prop:smc-resampling}
Conditional on the weighted sample $\{(x_t^{*,(i)},w_t^{(i)})\}_{i=1}^N$, let $x_t^{(1)},\dots,x_t^{(N)}$ be resampled independently from that empirical measure.
Then, for any integrable test function $h$,
\begin{equation}
  \E\!\left[
    \frac{1}{N}\sum_{j=1}^N h(x_t^{(j)})
    \,\bigm|\,
    x_t^{*,(1:N)},w_t^{(1:N)}
  \right]
  =
  \sum_{i=1}^N w_t^{(i)} h(x_t^{*,(i)}).
  \label{eq:smc-resampling-unbiased}
\end{equation}
\end{proposition}
\begin{proof}
Conditional on the current weighted sample, each $x_t^{(j)}$ has distribution $\sum_i w_t^{(i)}\delta_{x_t^{*,(i)}}$.
Therefore
\[
  \E[h(x_t^{(j)}) \mid x_t^{*,(1:N)},w_t^{(1:N)}]
  =
  \sum_{i=1}^N w_t^{(i)} h(x_t^{*,(i)}).
\]
Summing over $j=1,\dots,N$ and dividing by $N$ gives \eqref{eq:smc-resampling-unbiased}.%
\end{proof}

Resampling improves the distribution of future weights, but it is not free.
Duplicating the current high-weight particles reduces weight variance now, yet it also reduces genealogical diversity and can lead to sample impoverishment if resampling is too frequent.
That is why practical filters usually resample only when the ESS in \eqref{eq:smc-ess} falls below a threshold such as $N/2$.
For discussion of this tradeoff and of simple-random versus residual resampling schemes, see \citet[\S\S3.4.3--3.4.4]{liu_2002_monte_carlo}.

\begin{algorithm}[H]
  \caption{Adaptive bootstrap particle filter}
  \label{alg:bootstrap-particle-filter}
  \begin{algorithmic}[1]
    \Require Initial law $p_0$, transition density $f_\theta$, likelihood $g_\theta$, observations $y_{1:T}$, number of particles $N$, ESS threshold factor $c\in(0,1]$
    \For{$i=1,2,\dots,N$}
      \State Draw $x_0^{(i)} \sim p_0$ and set $w_0^{(i)} \gets 1/N$.
    \EndFor
    \For{$t=1,2,\dots,T$}
      \For{$i=1,2,\dots,N$}
        \State Draw $x_t^{*,(i)} \sim f_\theta(\cdot \mid x_{t-1}^{(i)})$.
        \State $\tilde w_t^{(i)} \gets w_{t-1}^{(i)} g_\theta(y_t \mid x_t^{*,(i)})$.
      \EndFor
      \State $\hat z_t \gets \sum_{i=1}^N \tilde w_t^{(i)}$.
      \For{$i=1,2,\dots,N$}
        \State $w_t^{(i)} \gets \tilde w_t^{(i)} / \hat z_t$.
      \EndFor
      \State $\mathrm{ESS}_t \gets 1 / \sum_{i=1}^N (w_t^{(i)})^2$.
      \If{$\mathrm{ESS}_t < cN$}
        \State Sample ancestor indices $A_t^{(1)},\dots,A_t^{(N)}$ with probabilities $w_t^{(1)},\dots,w_t^{(N)}$.
        \For{$i=1,2,\dots,N$}
          \State $x_t^{(i)} \gets x_t^{*,(A_t^{(i)})}$ and $w_t^{(i)} \gets 1/N$.
        \EndFor
      \Else
        \For{$i=1,2,\dots,N$}
          \State $x_t^{(i)} \gets x_t^{*,(i)}$.
        \EndFor
      \EndIf
    \EndFor
    \State \Return $\{(x_t^{(i)},w_t^{(i)})\}_{i=1}^N$ for $t=0,\dots,T$ and $\hat\ell_T(\theta)=\sum_{t=1}^T \log \hat z_t$
  \end{algorithmic}
\end{algorithm}

If the resampling step is carried out at every observation time, \Cref{alg:bootstrap-particle-filter} reduces to the classical sample-importance-resampling bootstrap filter.
More sophisticated resampling rules such as systematic, stratified, or residual resampling reduce the extra Monte Carlo variation introduced by multinomial resampling, but they play the same probabilistic role:
replace one highly skewed weighted cloud by another representation of essentially the same empirical law whose future weights will be less degenerate.

\subsubsection{Likelihood increments and parameter learning}
\label{subsubsec:smc-likelihood}

The filtering recursion also factorizes the marginal likelihood.
If
\begin{equation}
  z_t \coloneqq p_\theta(y_t \mid y_{1:t-1}),
  \label{eq:smc-likelihood-increment}
\end{equation}
then
\begin{equation}
  p_\theta(y_{1:T})
  =
  \prod_{t=1}^T z_t,
  \qquad
  \log p_\theta(y_{1:T})
  =
  \sum_{t=1}^T \log z_t.
  \label{eq:smc-likelihood-factorization}
\end{equation}
By \eqref{eq:smc-update}, each increment is a predictive expectation,
\begin{equation}
  z_t
  =
  \int g_\theta(y_t \mid x_t)\,p_\theta(x_t \mid y_{1:t-1})\,dx_t.
  \label{eq:smc-likelihood-predictive}
\end{equation}
The particle approximation therefore gives the estimator
\begin{equation}
  \hat z_t
  \coloneqq
  \sum_{i=1}^N \tilde w_t^{(i)}
  =
  \sum_{i=1}^N w_{t-1}^{(i)} g_\theta(y_t \mid x_t^{*,(i)}).
  \label{eq:smc-likelihood-estimator}
\end{equation}
When the particles entering time $t$ have already been reset to equal weights, this simplifies to
\begin{equation}
  \hat z_t
  =
  \frac{1}{N}\sum_{i=1}^N g_\theta(y_t \mid x_t^{*,(i)}).
  \label{eq:smc-likelihood-estimator-equal-weights}
\end{equation}
Combining \eqref{eq:smc-likelihood-factorization} and \eqref{eq:smc-likelihood-estimator} yields the natural particle approximation
\begin{equation}
  \hat p_\theta(y_{1:T})
  \coloneqq
  \prod_{t=1}^T \hat z_t,
  \qquad
  \hat\ell_T(\theta)
  \coloneqq
  \sum_{t=1}^T \log \hat z_t.
  \label{eq:smc-loglik}
\end{equation}
Even though the logarithm introduces bias by Jensen's inequality, the decomposition is still extremely useful because each term depends only on the current particle system and the newest observation.
This is exactly the structure one wants for online learning.
For this likelihood-increment factorization and its particle approximation, see \citet[\S2.6.3]{liu_2002_monte_carlo}.

The connection back to the latent-variable viewpoint from \Cref{subsec:expectation-maximization} is direct.
In a linear-Gaussian model, Kalman smoothing supplies exact conditional moments and EM updates the parameters analytically.
In a general nonlinear or non-Gaussian model, the same latent quantities are no longer available in closed form, but SMC supplies sequential approximations to the filtered laws and to the likelihood increments \eqref{eq:smc-likelihood-increment}.
Those approximations can be used for approximate EM, for likelihood-based tuning, or for stochastic-gradient parameter updates as new data arrive.
These parameter-learning uses of SMC are reviewed in \citet[\S\S3.4.3--3.4.4]{liu_2002_monte_carlo}.
The same complete-data logic from \eqref{eq:smc-em-sufficient-statistics} still applies:
if the M-step would be easy given latent trajectories, then particle smoothing or weighted ancestral trajectories provide approximations to the missing sufficient statistics, and one maximizes the resulting approximate $Q$-function.
For Bayesian parameter inference, one typically uses the likelihood estimate $\hat p_\theta(y_{1:T})$ rather than the log estimate \eqref{eq:smc-loglik}.
If the particle filter is constructed so that $\hat p_\theta(y_{1:T})$ is a nonnegative unbiased estimator of the true likelihood, then it can be used inside a pseudo-marginal Metropolis--Hastings update for $\theta$.
In that case the particle filter numerically integrates out the latent states while a separate sampler moves through parameter space.
Without that nonnegative unbiasedness property, the same substitution should be viewed only as an approximate plug-in.

Two practical cautions matter throughout.
First, state-space models can have identifiability symmetries, so parameter learning often needs constraints, regularization, or strong prior information to distinguish dynamically equivalent parametrizations.
Second, better approximation costs more computation:
more particles reduce Monte Carlo error, and more aggressive resampling stabilizes the ESS, but both also increase runtime and can reduce genealogical diversity.
In practice the first quantities to monitor are the ESS trajectory and the incremental log-likelihood estimates $\log \hat z_t$.
If the ESS repeatedly collapses or the likelihood estimate becomes erratic, the filter usually needs a better proposal, more particles, or a better constrained model.
